\section{Second Pass: Algebra as Lawful Manipulation}
\chapterline

In elementary high school algebra, the rules of arithmetic feel as natural as gravity. You freely add terms, rearrange factors, and divide both sides by whatever is in your way. 

Abstract algebra begins when you realize that those familiar manipulations were not laws of the universe, but consequences of specific axioms operating on specific number systems. When you study symmetries of geometric shapes, matrices, permutations, or cryptographic transformations, commutativity and division vanish. You must learn to navigate worlds governed strictly by declared laws.

\begin{historybox}
For centuries, algebra meant the art of solving polynomial equations with numbers. In 1930, a young Dutch mathematician named Bartel Leendert van der Waerden published \emph{Moderne Algebra}, based on the revolutionary lectures of Emmy Noether and Emil Artin at G\"ottingen and Hamburg. The book transformed the mathematical world. Instead of calculating explicit formulas for roots or computing gigantic matrix determinants, van der Waerden showed that algebra is the study of abstract axiomatic structures: groups, rings, fields, and ideals.

This structural revolution reached its zenith in the late twentieth century with the \emph{Classification of Finite Simple Groups}. Spanning several decades, thousands of pages, and hundreds of mathematicians (including Daniel Gorenstein, John Thompson, and Michael Aschbacher), this monumental achievement classified every atomic finite building block of symmetry into alternating groups, groups of Lie type, and 26 mysterious ``sporadic'' groups culminating in the Monster group.
\end{historybox}

\subsection{Why You Cannot ``Divide'': Left versus Right Inverses}

In a general group $(G, \cdot)$, the operation is associative, an identity element $e$ exists, and every element $g$ possesses an inverse $g^{-1}$ satisfying $g \cdot g^{-1} = g^{-1} \cdot g = e$. But $G$ is generally **non-commutative**: $a \cdot b$ and $b \cdot a$ can be completely different objects.

Consider the symmetries of a square (the dihedral group $D_4$). Let $r$ be a $90^\circ$ clockwise rotation, and let $s$ be a horizontal reflection. If you rotate first and reflect second, the result is $s \circ r$. If you reflect first and rotate second, the result is $r \circ s$. In fact,
\[
s \circ r = r^{-1} \circ s = r^3 \circ s \;\neq\; r \circ s.
\]
Because order matters, the concept of ``dividing by $a$'' is meaningless. Which side would you divide from? 

To solve $a \cdot x = b$, you must multiply by $a^{-1}$ on the **left**:
\[
a^{-1}(ax) = a^{-1}b \implies (a^{-1}a)x = a^{-1}b \implies ex = a^{-1}b \implies x = a^{-1}b.
\]
If you had multiplied on the right, you would obtain $axa^{-1} = ba^{-1}$, which does not isolate $x$ at all. In abstract algebra, precision of side is non-negotiable.

\subsection{Controversies, False Starts, and Gordan's ``Theology''}

The shift toward abstract axiomatic algebra encountered fierce resistance and pitfalls:

\begin{itemize}
  \item \textbf{Gordan's Cry of ``Theology!''}: In the late nineteenth century, Paul Gordan was known as the ``King of Invariants'' because of his extraordinary ability to calculate explicit generating systems for polynomial invariants. In 1890, David Hilbert proved the finite generation of invariants non-constructively using ideal theory and the Hilbert Basis Theorem, without producing a single explicit formula. Gordan famously cried: ``Das ist nicht Mathematik. Das ist Theologie!'' (This is not mathematics; it is theology!). Yet within a decade, Gordan embraced Hilbert's and Noether's structural approach, conceding that abstract axioms revealed truths that calculation could never touch.
  \item \textbf{The Phantom Zero-Divisor Trap}: In high school arithmetic, $ab = 0 \implies a = 0 \text{ or } b = 0$. In abstract rings, this fails routinely. In $\Z_6$, we have $2 \cdot 3 \equiv 0 \pmod 6$, even though neither $2$ nor $3$ is zero! Elements that multiply to zero are \emph{zero divisors}. If you naively cancel $a$ from $ab = ac$, you may divide by a zero divisor, creating false proofs like $1 = 2$. Cancellation is lawful only in integral domains and division rings.
  \item \textbf{The Non-Normal Quotient Collapse}: Students often assume that for any subgroup $H \le G$, the set of cosets $G/H$ forms a quotient group. But if $H$ is not normal, coset multiplication is not well-defined! Choosing different elements from the same cosets yields products in conflicting cosets, causing the entire algebraic structure to disintegrate.
\end{itemize}

\subsection{The Kernel: Measuring Lost Information}

When a homomorphism $\phi: G \to H$ maps one group into another, it respects the group multiplication:
\[
\phi(a \cdot b) = \phi(a) \cdot \phi(b).
\]
The elements that get crushed down to the identity $e_H \in H$ form the **kernel**:
\[
\ker(\phi) = \{g \in G : \phi(g) = e_H\}.
\]
The kernel is never just an arbitrary collection of elements; it is a subgroup of $G$ with a very special geometric and algebraic symmetry called **normality**.

\begin{center}
\begin{tikzpicture}[scale=1.1, >=Stealth]
  \tikzset{
    groupbox/.style={draw=sourceblue!80!black, fill=sourceblue!5, very thick, rounded corners=4pt},
    cosetbox/.style={draw=sourceblue!70!black, fill=sourceblue!15, thick, rounded corners=3pt},
    kernbox/.style={draw=structurepurple!90!black, fill=structurepurple!20, very thick, rounded corners=3pt},
    pointdot/.style={circle, fill=sourceblue!90!black, inner sep=1.8pt},
    targbox/.style={draw=targetgreen!80!black, fill=targetgreen!6, very thick, rounded corners=4pt},
    targpoint/.style={circle, fill=targetgreen!90!black, inner sep=2.5pt},
    homarrow/.style={->, thick, shorten >=3pt, shorten <=3pt}
  }

  % Group G (Domain)
  \node[font=\footnotesize\scshape, text=sourceblue!80!black] at (2.0, 4.3) {Domain Group $G$ (Coset Partition)};
  \draw[groupbox] (0, 0) rectangle (4.0, 3.8);

  % Coset 1: Kernel N = ker(\phi)
  \draw[kernbox] (0.3, 2.7) rectangle (3.7, 3.5);
  \node[pointdot, fill=structurepurple!90!black] at (1.0, 3.1) {};
  \node[pointdot, fill=structurepurple!90!black] at (2.0, 3.1) {};
  \node[right=2pt, font=\small\bfseries, text=structurepurple!90!black] at (2.3, 3.1) {$N = \ker(\phi)$};
  \node[above=1pt, font=\tiny] at (1.0, 3.1) {$e_G$};

  % Coset 2: g1 N
  \draw[cosetbox] (0.3, 1.8) rectangle (3.7, 2.5);
  \node[pointdot] at (1.2, 2.15) {};
  \node[pointdot] at (2.4, 2.15) {};
  \node[right=2pt, font=\small, text=sourceblue!80!black] at (2.6, 2.15) {$g_1 N$};

  % Coset 3: g2 N
  \draw[cosetbox] (0.3, 0.9) rectangle (3.7, 1.6);
  \node[pointdot] at (1.5, 1.25) {};
  \node[pointdot] at (2.7, 1.25) {};
  \node[right=2pt, font=\small, text=sourceblue!80!black] at (2.9, 1.25) {$g_2 N$};

  % Coset 4: g3 N
  \draw[cosetbox] (0.3, 0.2) rectangle (3.7, 0.75);
  \node[right=2pt, font=\small, text=sourceblue!80!black] at (1.8, 0.48) {$g_3 N$};

  % Target Group H
  \node[font=\footnotesize\scshape, text=targetgreen!80!black] at (8.0, 4.3) {Target Group $H$};
  \draw[targbox] (6.5, 0) rectangle (9.5, 3.8);
  \node[font=\scriptsize\scshape, text=targetgreen!80!black] at (8.0, 3.5) {$\operatorname{im}(\phi) \subseteq H$};

  % Points in H
  \node[targpoint, fill=structurepurple!90!black, label={right:\small $e_H = \phi(N)$}] (h0) at (7.2, 3.1) {};
  \node[targpoint, label={right:\small $\phi(g_1)$}] (h1) at (7.2, 2.15) {};
  \node[targpoint, label={right:\small $\phi(g_2)$}] (h2) at (7.2, 1.25) {};
  \node[targpoint, label={right:\small $\phi(g_3)$}] (h3) at (7.2, 0.48) {};

  % Homomorphism arrows between cosets and image points
  \draw[homarrow, draw=structurepurple!90!black, thick] (3.8, 3.1) -- (7.0, 3.1);
  \draw[homarrow, draw=sourceblue!80!black] (3.8, 2.15) -- (7.0, 2.15);
  \draw[homarrow, draw=sourceblue!80!black] (3.8, 1.25) -- (7.0, 1.25);
  \draw[homarrow, draw=sourceblue!80!black] (3.8, 0.48) -- (7.0, 0.48);

  \node[above, font=\footnotesize, text=gray!70!black] at (5.4, 3.2) {All of $N$ collapses to $e_H$};
  \node[below, font=\footnotesize, text=gray!70!black] at (5.4, 1.7) {Each coset maps to 1 point};
\end{tikzpicture}
\end{center}

\subsection{Normal Subgroups: Why Coset Multiplication Must Work}

If we want to form a quotient group $G/N$ whose elements are the cosets $gN$, we need a rule to multiply them:
\[
(aN) \cdot (bN) \stackrel{?}{=} (ab)N.
\]
This looks deceptively simple, but there is an immediate danger: this definition refers to specific representatives $a$ and $b$. What if we picked different elements from the exact same cosets, say $a' \in aN$ and $b' \in bN$? Would the product $(a'b')N$ land in the identical coset $(ab)N$?

In general, **no**! Coset multiplication is well-defined if and only if the subgroup $N$ is **normal**:
\[
gNg^{-1} = N \quad \text{for all } g \in G.
\]
\begin{formalbox}
A subgroup $N \le G$ is normal (written $N \triangleleft G$) if it is invariant under conjugation by every element of $G$. In other words, for every $n \in N$ and $g \in G$, the element $gng^{-1}$ remains safely inside $N$.
\end{formalbox}

Every kernel is automatically a normal subgroup. If $n \in \ker(\phi)$ and $g \in G$:
\[
\phi(gng^{-1}) = \phi(g)\phi(n)\phi(g)^{-1} = \phi(g) e_H \phi(g)^{-1} = \phi(g)\phi(g)^{-1} = e_H.
\]
Therefore $gng^{-1} \in \ker(\phi)$. 

The First Isomorphism Theorem for groups now crystallizes:
\[
G/\ker(\phi) \;\cong\; \operatorname{im}(\phi).
\]
The quotient group $G/\ker(\phi)$ is an exact algebraic mirror of the image group inside $H$.

\subsection{What Abstract Algebra Unlocks---and Where It Hits a Wall}

The structural algebra of Noether and van der Waerden revolutionized modern science:
\begin{enumerate}
  \item \textbf{Galois Theory and Insolvability}: Algebra solves centuries-old geometric questions---proving that no algebraic formula in radicals can solve general quintic polynomials, and that angle trisection and circle squaring are impossible with straightedge and compass.
  \item \textbf{The Symmetry Language of Physics}: In quantum field theory, fundamental particles are classified as irreducible representations of Lie groups: quarks, leptons, and gauge bosons correspond to the representations of $\mathrm{SU}(3) \times \mathrm{SU}(2) \times \mathrm{U}(1)$.
  \item \textbf{Modern Cryptography}: Public-key cryptosystems (RSA, Diffie--Hellman, and Elliptic Curve Cryptography) rely entirely on the asymmetry of algebraic group operations (discrete logarithms and modular exponentiation).
\end{enumerate}

\paragraph{Honest Downsides and Limitations.}
Yet algebra has clear limitations:
\begin{itemize}
  \item \textbf{Algorithmic Undecidability (The Word Problem)}: Although group axioms are simple, computing with them can be provably impossible. The Novikov--Boone theorem (1955) proved that there exists a finitely presented group for which no algorithm can determine whether an arbitrary word represents the identity element.
  \item \textbf{Blindness to Geometric Distance}: Purely algebraic rings and groups have no concept of metric proximity or geometric holes. A ring cannot tell whether its geometric spectrum is connected, compact, or curved without topological tools.
\end{itemize}

\paragraph{Where We Go Next in the Book.}
To understand continuity, shape, and nearness without being trapped by formulas or arithmetic, we must study spaces directly. In Chapter~37 (\emph{Second Pass: Topology as the Study of Stable Nearness}), we throw away rulers and coordinate axes completely, defining continuity through the open set.

\subsection{Exercises}

\begin{enumerate}
  \item In any group $G$, prove the socks-and-shoes property: for all $a, b \in G$,
  \[
  (ab)^{-1} = b^{-1}a^{-1}.
  \]
  Explain using the metaphor of putting on socks and shoes why the order must reverse.
  \item Let $G$ be a group and let $H \le G$ be a subgroup of index $2$ (meaning there are exactly two cosets: $H$ and $G \setminus H$). Prove that $H$ must be a normal subgroup of $G$.
  \item Consider the sign homomorphism $\operatorname{sgn}: S_n \to \{+1, -1\}$ on the symmetric group of permutations, which sends even permutations to $+1$ and odd permutations to $-1$.
  \begin{enumerate}
    \item What is the kernel $\ker(\operatorname{sgn})$? (This is called the alternating group $A_n$.)
    \item What is the size $|A_n|$ relative to $|S_n| = n!$?
    \item Apply the First Isomorphism Theorem to deduce $S_n/A_n \cong \Z_2$.
  \end{enumerate}
\end{enumerate}

\newpage
